\section{Conclusion}
\label{sec:conslusion}

In this paper, we propose MoDiff, a principled framework for accelerating generative modeling. Our preliminary studies reveal the challenges in caching and PTQ methods. To address these, we introduce modulated quantization and error compensation, which reduce quantization error and mitigate error accumulation. We provide theoretical analyses demonstrating the effectiveness of our approach. Experimental results show that MoDiff significantly enhances activation quantization, enabling PTQ methods to operate at bit-widths as low as 3 bits without performance degradation. One limitation is that MoDiff reduces computation at the cost of increased memory usage. Additionally, we evaluate acceleration based on theoretical computational complexity rather than real-world hardware speedup. We leave hardware implementation and further memory optimizations of our MoDiff for future work.

\newpage

\section*{Acknowledgments}

Weizhi Gao, Zhichao Hou, and Dr. Xiaorui Liu are supported by the National Science Foundation (NSF) National AI Research Resource Pilot Award, Amazon Research Award, NCSU Data Science Academy Seed Grant Award, and NCSU Faculty Research and Professional Development Award. 

\section*{Impact Statement}
\textbf{Ethical Considerations.} This work introduces modulated diffusion models to accelerate the generation process of diffusion models. We do not foresee any significant ethical concerns associated with this approach.

\textbf{Societal Impact.} Enhancing the efficiency of diffusion models can facilitate the broader adoption of generative AI, benefiting both AI research and hardware development. This advancement has the potential to contribute to more efficient and accessible AI-driven solutions.
